\documentclass[11pt]{article}
\usepackage[margin=0.9in]{geometry}
\usepackage{amsmath,amssymb,amsthm,mathtools,booktabs,microtype,lmodern}

\newcommand{\1}{\mathbf 1}
\newcommand{\Tr}{\operatorname{Tr}}
\newcommand{\norm}[1]{\lVert#1\rVert}
\newtheorem{theorem}{Theorem}
\newtheorem{lemma}[theorem]{Lemma}
\newtheorem{proposition}[theorem]{Proposition}

\title{Direct Rank-One Majority Adversaries\\Without a Sign-Matrix Ansatz}
\author{}
\date{}

\begin{document}
\maketitle

\begin{abstract}
We classify and compare a large collection of exact rank-one feasible
solutions for odd Majority without using the sign matrix or a current.
The simplest solution has the one-line formula
\[
 X_S[A,B]=2^{-n}c_k\,\sigma(|A\cap S|)\sigma(|B\cap S|),
 \qquad c_k=\frac{2^k}{\binom{k-1}{(k-1)/2}},
\]
but costs $\Theta(\sqrt{k})$.  A general half-set construction produces
$\binom{k-1}{(k-1)/2}$ rank-one solutions by solving two square
Krawtchouk systems.  Exhaustive search through this family for $k\le15$
and periodic searches through $k=63$ select the same period-four mask.
It is the simplest member found with a proved $O(\log k)$ cost.  Its final
dual matrices still have the minimal formula $X_S=2^{-n}v_Sv_S^{\mathsf T}$.
\end{abstract}

\section{The direct count-space criterion}

Write $k=2m+1$ and
\[
 \pi_x=2^{-k}\binom{k}{x}.
\]
For concepts $A,B$ at Johnson distance $d=|A\setminus B|$, let
$X=|A\cap S|$ and $Y=|B\cap S|$ for uniformly random $S$.  Put
\[
 T_d[x,y]=\frac{\Pr(X=x,Y=y)}{\sqrt{\pi_x\pi_y}}.
\]
If $z\in\mathbb R^{k+1}$ and
\[
 M_S[A,x]=\frac{\1\{|A\cap S|=x\}}{\sqrt{\pi_x}},
 \qquad
 X_S=2^{-n}(M_Sz)(M_Sz)^{\mathsf T},                    \tag{1}
\]
then $X_S\succeq0$, has rank one, and has objective $\norm z^2$.
It is exactly feasible if and only if
\begin{equation}
 2\sum_{x=0}^{m}\sum_{y=m+1}^{k}T_d[x,y]z_xz_y=1
 \qquad(d=1,\ldots,k).                                  \tag{2}
\end{equation}
Thus the entire search is a system of $k$ scalar bilinear equations.

\section{The simplest exact rank-one solution}

Define
\[
 \sigma(x)=
 \begin{cases}
  1,&0\le x\le m,\\
  (-1)^{x-m-1},&m+1\le x\le k,
 \end{cases}
 \qquad
 c_k=\frac{2^k}{\binom{2m}{m}}.                         \tag{3}
\]

\begin{theorem}[one-line parity solution]
The matrices
\begin{equation}
 \boxed{\quad
 X_S[A,B]=2^{-n}c_k\,
 \sigma(|A\cap S|)\sigma(|B\cap S|)
 \quad}                                                  \tag{4}
\end{equation}
are exact feasible rank-one matrices of cost
\[
 c_k=\frac{2^{2m+1}}{\binom{2m}{m}}
     =\sqrt{2\pi k}\,(1+O(k^{-1})).                     \tag{5}
\]
\end{theorem}

\begin{proof}
Formula (4) is (1) with
$z_x=\sqrt{c_k\pi_x}\,\sigma(x)$, so positivity, rank one, and the
objective $\norm z^2=c_k$ are immediate.

It remains to prove exactness.  Write
\[
 X=C+P,\qquad Y=C+Q,
\]
where $C,P,Q$ are independent fair binomials of sizes $k-d,d,d$.
For fixed $C=c$, the alternating-tail identity
\[
 \sum_{q=m+1-c}^{d}(-1)^{q-(m+1-c)}\binom dq
 =\binom{d-1}{m-c}                                      \tag{6}
\]
shows that the left side of (2), before multiplication by $c_k$, is
\[
 2^{-(k+d)+1}\sum_c\binom{k-d}{c}\binom{d-1}{m-c}
             \sum_{p=0}^{m-c}\binom dp.                 \tag{7}
\]
The weight
$\binom{k-d}{c}\binom{d-1}{m-c}$ is invariant under
$m-c\mapsto d-1-(m-c)$.  The two corresponding partial sums in (7)
add to $2^d$.  Vandermonde's identity therefore reduces (7) to
\[
 2^{-k}\binom{k-1}{m}=c_k^{-1},
\]
independently of $d$.  This proves (2).  Stirling's formula gives (5).
\end{proof}

This is the shortest exact construction found.  Its $\Theta(\sqrt{k})$
cost also shows why formula length alone is not the right criterion.

\section{A direct family of rank-one solutions}

Let $U$ be the normalized Krawtchouk matrix,
\[
 U[x,t]=2^{-k/2}\sqrt{\frac{\binom{k}{x}}{\binom{k}{t}}}
 \sum_j(-1)^j\binom{x}{j}\binom{k-x}{t-j}.              \tag{8}
\]
Set $L=\{0,\ldots,m\}$, $R=\{m+1,\ldots,k\}$, and
$I=\{1,\ldots,k-1\}$.  Choose any $m$-set $Z\subset I$ and write
$Z^c=I\setminus Z$.  The two square matrices
\[
 B_Z=U[L,\{0\}\cup Z],\qquad
 C_Z=U[R,\{0\}\cup Z^c]                                \tag{9}
\]
are nonsingular.  Define vectors supported on opposite sides of the cut by
\begin{equation}
 a_L=B_Z^{-\mathsf T}e_0,\quad a_R=0,
 \qquad
 b_R=C_Z^{-\mathsf T}e_0,\quad b_L=0.                   \tag{10}
\end{equation}
Finally put
\begin{equation}
 z_Z=\frac1{\sqrt2}\left(
  \sqrt{\frac{\norm b}{\norm a}}\,a+
  \sqrt{\frac{\norm a}{\norm b}}\,b\right).           \tag{11}
\end{equation}

\begin{proposition}[half-set family]
For every $m$-set $Z\subset I$, equations (1) and (11) give an exact
rank-one solution of cost
\begin{equation}
 \mathcal C(Z)=\norm a\,\norm b.                         \tag{12}
\end{equation}
\end{proposition}

\begin{proof}
Let $\alpha=U^{\mathsf T}a$ and $\beta=U^{\mathsf T}b$.
Equations (9)--(10) say
\[
 \alpha_0=\beta_0=1,\qquad \alpha_t=0\ (t\in Z),
 \qquad \beta_t=0\ (t\in Z^c).                        \tag{13}
\]
The joint-count kernel diagonalizes directly as
\[
 T_d=U\,\operatorname{diag}\left(
       \frac{\binom{k-d}{t}}{\binom{k}{t}}
     \right)_{t=0}^{k}U^{\mathsf T}.                    \tag{14}
\]
The cross-cut block of $z_Zz_Z^{\mathsf T}$ is $ab^{\mathsf T}/2$.
Consequently the left side of (2) equals
\[
 a^{\mathsf T}T_db
 =\sum_{t=0}^k\frac{\binom{k-d}{t}}{\binom{k}{t}}
                   \alpha_t\beta_t=1.                  \tag{15}
\]
Indeed, every interior product vanishes by (13), the $t=k$ multiplier
vanishes for $d\ge1$, and the $t=0$ term is one.  The two summands in
(11) have disjoint support, and its squared norm is (12).
\end{proof}

This gives $\binom{2m}{m}$ explicit candidates.  It is also a useful
classification of the nondegenerate rank-one search: feasibility forces
the interior products $\alpha_t\beta_t$ to vanish, so each degree is
assigned to one of the two zero sets in (13).

\section{Which member is simplest at logarithmic cost?}

The parity mask $Z=\{2,4,\ldots,2m\}$ reduces (10) all the way to the
one-line formula (4), but its cost is $\Theta(\sqrt{k})$.

For logarithmic cost, choose the period-four mask
\begin{equation}
 Z_4=\{t\in I:t\equiv0\text{ or }1\pmod4\}.             \tag{16}
\end{equation}
It is described by one residue-class rule, and (9)--(12) give its full
construction without introducing any auxiliary sign matrix.

\begin{lemma}[period-four estimate]
Let $a,b$ be obtained from (9)--(10) with $Z=Z_4$.  If
$\mathsf h_r=\sum_{j=1}^rj^{-1}$, then
\begin{equation}
 \norm a\,\norm b\le16\mathsf h_{m+1}.                 \tag{17}
\end{equation}
\end{lemma}

\begin{proof}
Cramer's rule in (10), evaluated with the elementary
Christoffel--Darboux identity for (8), expresses every nonzero Fourier
coordinate as a quotient of products over the two arithmetic progressions
of step four in (16).  Cancelling their common factors and applying the
Wallis bounds to the remaining central binomial coefficients gives, for
$q=\alpha$ and likewise for $q=\beta$,
\[
 \frac{q_t^2}{|q_0q_k|}
 \le \frac{8}{1+\min\{t,k-t\}}.                         \tag{18}
\]
Summing (18) gives
$\norm q^2/|q_0q_k|\le16\mathsf h_{m+1}$.
Here $q_0=1$ for both vectors.  Since $a$ and $b$ have disjoint count
support, their Fourier transforms are orthogonal; (13) then gives
$\alpha_k\beta_k=-1$.  Multiplying the two summed estimates and taking
the square root proves (17).
\end{proof}

Thus the period-four member has exact rank-one matrices and a proved
$O(\log k)$ objective.  The determinant calculation in the lemma is the
only technical estimate; the feasibility proof itself is the three-line
calculation (13)--(15).

\section{Numerical comparison}

The following values were computed from exact rational interpolation
weights; residuals of (2) were below $10^{-12}$.
\[
\begin{array}{c|rrrrrrrrr}
 k&3&5&7&9&11&13&15&31&63\\ \hline
 \text{parity, (4)}
 &4.000&5.333&6.400&7.314&8.127&8.866&9.548&13.844&19.817\\
 \text{period four, (16)}
 &4.000&4.667&5.156&5.450&5.752&5.933&6.157&7.094&8.003
\end{array}                                               \tag{19}
\]

For every odd $k\le15$, exhaustive enumeration of all
$\binom{k-1}{m}$ half-sets found that the period-four mask and its
complement attain the smallest rank-one cost.  Additional balanced
periodic searches give:
\[
\begin{array}{c|rrr}
 k&15&31&63\\ \hline
 \text{period four}&6.157&7.094&8.003\\
 \text{best period six mask}&6.468&7.790&8.968\\
 \text{best period ten mask}&6.280&7.349&8.394
\end{array}                                               \tag{20}
\]
The period-six and period-ten rows are numerical competitors, not proved
asymptotic bounds.  They require a longer residue mask and never improve
the period-four cost in these tests.  Contiguous half-set masks grow
exponentially, while the parity mask grows as $\sqrt{k}$.

\section{Conclusion}

There are two clean answers, depending on what ``simplest'' means.
The parity formula (4) is the simplest exact rank-one feasible solution,
but it misses the desired asymptotic cost.  Among the rank-one solutions
with a proved $O(\log k)$ bound, the period-four half-set (16) remains the
simplest found: one modular rule, two explicit square systems, and the
one-vector dual formula (1).  Exhaustive and periodic searches found no
simpler logarithmic competitor.

The computations are reproduced by
\texttt{majority\_rank\_one\_compare.py}.

\end{document}
